\appendix
\section{Formal Specification}
\label{app:formal}
% ================================================================================

This appendix provides the formal artefacts accompanying the
Bridge Contract specification. The artefacts are also available as
standalone files in the accompanying \texttt{formal-spec/} directory:
\texttt{BridgeContract.tla}, the TLC configurations, a mutation script
and a \texttt{README.md} with run instructions and the recorded results.
The product machine of Section~\ref{sec:product-model} is the normative
state model stated here, and the TLA+ module models it (Appendix~\ref{app:tla-extract}).

\subsection{Product model}
\label{app:product-model}

The formal state space is the product
\[
\mathcal{S} = \textsc{Attachment} \times \textsc{Human} \times \textsc{Lifecycle}
\]
where
\begin{itemize}[leftmargin=2em]
\item $\textsc{Attachment} = \{\textsc{Attached},
\textsc{Detached}, \textsc{Resuming}\}$;
\item $\textsc{Human} = \{\textsc{None}, \textsc{Pending}\}$;
\item $\textsc{Lifecycle} = \{\textsc{Active}, \textsc{Terminated}\}$.
\end{itemize}

The product space is Attachment $\times$ Human $\times$ Lifecycle. Every
active cell is reachable; the terminated cells are folded into
one, $\bot$:

\footnotesize
\begin{longtable}{l l l p{0.40\columnwidth}}
\caption{Product cells and their derived named-state view.}\label{tab:cells}\\
\toprule
\textbf{Attachment} & \textbf{Human} & \textbf{Lifecycle} & \textbf{Derived view of named states} \\
\midrule
Attached & None & Active & ATTACHED \\
Detached & None & Active & DETACHED \\
Attached & Pending & Active & PENDING\_HUMAN \\
Detached & Pending & Active & PENDING\_HUMAN \\
Resuming & None & Active & RESUMING \\
Resuming & Pending & Active & RESUMING \\
any & any & Terminated & TERMINATED ($\bot$) \\
\bottomrule
\end{longtable}

The fold of $(\textsc{Att},\textsc{Pend})$ and $(\textsc{Det},\textsc{Pend})$
into one name is lossy, which is why the named-state view is not
normative (Section~\ref{sec:mece-status}).

\subsection{TLA+ model and TLC results}
\label{app:tla-extract}

\texttt{formal-spec/BridgeContract.tla} models the product machine with
the client set (I15), challenge-response re-attach with a single-use nonce
and a prefix check (I4, I9, I11), narrow-only authority (I3, I10), durable
append and restart recovery (I12), the tool queue (I13) and answer binding
(I14), quarantine on a chain failure (rule~Q), the cumulative detachment
cap $T_{\text{total}}$, the detached-call policy (rule~D, calls outside an
escalation), the forwarding acknowledgement and the key lifecycle
(I11-REVOKE). The key set $K$ keeps revoked keys, marked as revoked
(\texttt{rev}) and, for a revocation for compromise, as hard
(\texttt{hard}); the actions are a soft or hard revocation (KEY\_REVOKED; a
hard one also detaches a connected client and abandons a re-attach that
client is verifying), \texttt{OperatorAddKey} (OPERATOR\_ADD\_KEY) and
\texttt{OperatorEnd} (OPERATOR\_TERMINATION, also for a key-exhausted
session). A session in which every registered key is revoked is
key\_exhausted, a side predicate rather than a cell: it keeps its cell and
serves no client action, but its timers keep running, a restart treats it like
any other session and a lost transport detaches it, so a detached
key-exhausted session ends after $T_{\text{detach}}$ unless
\texttt{OperatorAddKey} adds a key; \texttt{OperatorEnd} ends it at once. A revoked key
is refused at re-attach, which still spends the nonce it names. A revocation of the
key being verified, hard or soft, abandons the verification; both reference
bridges abandon the re-attach in the same way (T128, V111). Invariants carry
the paper's names (\texttt{IRev}: a key revoked for compromise is neither
connected nor being verified); step properties
($[\,][\,]_{\mathit{vars}}$ formulas) state I3, I4, I5--I7, I9, I10, I11,
I11-REVOKE, I14, I15 and Rule~D over single transitions, among them that
$K$, \texttt{rev} and \texttt{hard} only grow, that a joining client is not
revoked, that a key is registered only through a connected unrevoked client or
by the operator on a key-exhausted session, and that a key-exhausted session
changes only by time, a restart or a lost transport, the operator's end or a
new key. The event log is represented by its
length and the state it determines, a transition appends its events and
changes the state in one committed step, and a crash between the append and the in-memory change is
equivalent to the same step followed by a restart, which rebuilds the
state from the log. Hashes are abstract. A re-attach is presented as a key, a signature, a
claimed nonce and a prefix verdict; the signature is a record that verifies
only for the key, nonce and verdict it names (Appendix~\ref{app:attacker}),
and only an authentic request with a live unspent nonce and a matching prefix
is accepted. The nonce book keeps every nonce issued with its own countdown;
at most a bounded number are live and a challenge is refused when the book is
full (I11 (f)); an unauthenticated failure is recorded only while fewer than a
bounded number were recorded since the last request whose signature verified
(I11 (g)).

SANY parses the module and TLC reports no error in every configuration
(Table~\ref{tab:tlc}). The runs used the TLC build
\texttt{2026.11.19.180144} (the revision of the TLA+ tools jar, \texttt{bc1b506}, not of this repository) from the TLA+ tools release
(sha256 \texttt{fdffef24\ldots 6713}; the banner is a build
number, not a release number such as 2.19), OpenJDK 21 in the container
\texttt{eclipse-temurin:21-jre}, one core and one worker per run. The image digest is
\texttt{sha256:d77e090a\ldots de0d}; the
logs are in \texttt{docs/logs/tlc-rerun-2026-11-25/} of the repository. No run
was resumed from a checkpoint, so every count is exact; the configurations ran one
at a time. The counts are sizes, not values a rerun must reproduce
to the last digit.

\begin{table}[htbp]
\centering
\caption{TLC results for \texttt{BridgeContract.tla}. The safety rows count
classes of states under a view that drops what cannot affect the checked
properties (see the text); the fair row runs without it.}
\label{tab:tlc}
\small
\begin{tabularx}{\columnwidth}{X l r r r}
\toprule
\textbf{Configuration} & \textbf{Checks} & \textbf{Generated} & \textbf{Distinct} & \textbf{Depth} \\
\midrule
2 sessions, \texttt{MaxEvents}=2, no attacker & safety & 13,591,683 & 2,883,227 & 21 \\
1 session, \texttt{MaxEvents}=7, no attacker & safety & 42,459,668 & 7,896,325 & 21 \\
1 session, \texttt{MaxEvents}=5, attacker, \texttt{ObsMax}=2 & safety & 17,159,349 & 444,499 & 18 \\
1 session, \texttt{MaxEvents}=5, fair & L1, L1Res, L2, cap & 6,358,094 & 1,391,382 & 20 \\
1 session, \texttt{MaxEvents}=5, policy \texttt{queue} & safety & 2,989,885 & 661,356 & 18 \\
1 session, \texttt{MaxEvents}=6, policy \texttt{forward} & safety & 11,321,712 & 2,259,734 & 20 \\
1 session, \texttt{MaxEvents}=5, $T_{\text{human}} > T_{\text{detach}}$ & safety & 3,714,583 & 821,343 & 18 \\
1 session, \texttt{MaxEvents}=5, one permission, up to two live nonces and two recorded failures & safety & 2,668,635 & 425,861 & 18 \\
\bottomrule
\end{tabularx}
\end{table}

\paragraph{What each configuration reaches.} A configuration that
reaches few cells or fires few actions checks little, whatever its state
count. Table~\ref{tab:tlc-reach} lists, for each configuration, the cells
that no state is in and the actions that generate no new state
(\texttt{formal-spec/tlc\_coverage.py}; the cells are the active cells,
$\bot$, quarantine and the key\_exhausted side predicate; the fair configuration is probed without fairness,
which does not change which states exist). The two-session run, with an event bound of 2,
never reaches a cell with $a=\textsc{Res}$ and never fires
\texttt{Disconnect}, \texttt{Resolve}, \texttt{Submit}, \texttt{Drain},
\texttt{Forward}, \texttt{OperatorAddKey}, \texttt{OperatorEnd} or the
re-attach completion: a detachment appends two
events and so exceeds the bound. What it checks is the interaction of two
sessions through the clock and a restart, over escalation, narrowing and
termination, and key revocation (\texttt{Revoke} fires and the key\_exhausted
state is reached, but the operator's add-key and end from it are not); the
detach, resume, drain and key-exhaustion exits are checked by the one-session
runs, which reach every cell and the key\_exhausted state. \texttt{SubmitAway} never fires under the policy
\texttt{refuse}, because a refused call changes nothing. \texttt{Deliver},
the delivery acknowledgement, fires only in the run with the larger event
bound; the other runs do not exercise it, and the acknowledgement mutants
are run on a deeper, narrower configuration for that reason. The adversarial
form of \texttt{Present} fires only in the attacker configuration, as the
attacker is off in the others, and is left out of the table for that reason.
The probes ran one configuration at a time, each as one uninterrupted run, and
their results are in \texttt{docs/logs/tlc-coverage-2026-11-24/}.

\begin{table}[htbp]
\centering
\caption{Cells not reached and actions that never fire, per TLC
configuration.}
\label{tab:tlc-reach}
\small
\begin{tabularx}{\columnwidth}{X l X}
\toprule
\textbf{Configuration} & \textbf{Cells not reached} & \textbf{Actions that never fire} \\
\midrule
2 sessions, \texttt{MaxEvents}=2, no attacker & $(\textsc{Res},\textsc{None})$,
$(\textsc{Res},\textsc{Pend})$ & \texttt{Answer}, \texttt{CompleteResume},
\texttt{Deliver}, \texttt{Disconnect}, \texttt{Drain},
\texttt{ExpireCall}, \texttt{Extend}, \texttt{Forward},
\texttt{OperatorAddKey}, \texttt{OperatorEnd}, \texttt{Resolve}, \texttt{Submit},
\texttt{SubmitAway} \\
1 session, \texttt{MaxEvents}=7, no attacker & none & \texttt{SubmitAway} \\
1 session, \texttt{MaxEvents}=5, attacker, \texttt{ObsMax}=2 & quarantine &
\texttt{Corrupt}, \texttt{Deliver}, \texttt{SubmitAway} \\
1 session, \texttt{MaxEvents}=5, fair & none & \texttt{Deliver},
\texttt{SubmitAway} \\
1 session, \texttt{MaxEvents}=5, policy \texttt{queue} & none &
\texttt{Deliver} \\
1 session, \texttt{MaxEvents}=6, policy \texttt{forward} & none &
\texttt{Deliver} \\
1 session, \texttt{MaxEvents}=5, $T_{\text{human}} > T_{\text{detach}}$ &
none & \texttt{Deliver}, \texttt{SubmitAway} \\
1 session, \texttt{MaxEvents}=5, one permission, up to two live nonces and two recorded failures &
quarantine & \texttt{Corrupt}, \texttt{Deliver}, \texttt{SubmitAway} \\
\bottomrule
\end{tabularx}
\end{table}

\paragraph{State-space reduction and what it costs.} Adding the attacker to
every configuration made the deep run too large (it had not finished after
five hours at over $10^8$ distinct states), so the model has two switches and
a view. The attacker is on only in the attacker configuration (one session,
one permission, one call key, no log tampering, \texttt{MaxEvents}=5); in
the others it is off, so the configuration with two sessions has no attacker. The
attacker remembers at most \texttt{ObsMax}=2 honest signatures. Hence the
attacker checks (\texttt{NoForgery}, replay, forgery with its own keys) run
only in that configuration; its interleaving with the tool queue, draining
and forwarding, and with log corruption and quarantine (the cell is not
reached there), is not checked; and an attack that combines three or more
distinct remembered signatures is not explored. The nonce-book and
failure-record invariants stay in every configuration. The safety
configurations use a view that drops the last presentation, the spent-nonce
set and the countdowns of dead nonces, the approval countdown when it cannot
be read, and all but the outcome of a terminated session; step properties are
evaluated on every transition. The argument that the view is sound is in the
README of the specification; the queue, forward and long human timeout configurations use it too.

To check that the verdict is not vacuous, mutants of the module are
defined (for example, for the key lifecycle, a revoked key accepted, a
revocation that removes the key from $K$ or accepts a key outside it, a hard
revocation that leaves the client connected or does not abandon a
verification, a revoked client that registers a key, a key-exhausted session
that serves client operations, whose timers stop or whose timeout does not fire,
OPERATOR\_ADD\_KEY on a session that is not key-exhausted or that clears the
revocation markers; for the attacker model, a client's signing key in the attacker's hands, a signature that verifies under any key or that is not bound to its nonce or to the prefix verdict; for I11 (f) and (g), a challenge that ignores the live-nonce limit or evicts a live nonce, unauthenticated failures that are all recorded, a verifying request dropped past the limit and an authentic request that does not restore the allowance; and further an expired or spent nonce accepted, the prefix not
checked, a partial departure that detaches, a drain that ignores the answer
hash, a widening applied, a chain failure that terminates instead of
quarantining, a cap that a restart or a disconnect resets or that counts
time attached, a delivered call forwarded again, a forward that ignores
the operator's allowlist). TLC rejects every mutant in the set, run on one session with \texttt{MaxEvents}=6 under the same view (the attacker mutants with the attacker on; the last line of
\texttt{docs/logs/tlc-rerun-2026-11-25/mutants.log} reports all killed; each yields a violation
of an invariant or of the step properties), and removing the fairness of
the timeout action makes L1 and L2 fail on this model. Findings from the
model: L1 is verified in its per-episode form (a detached session leaves
$\textsc{Det}$, terminates or is quarantined), and it showed that a
quarantined session is the one cell besides $\bot$ in which a session may
stay for ever. An earlier form of the model froze a key-exhausted session,
so that no timer ended it, and L1 and L2 had to name it as an exception;
the model now keeps its timers running and L1 and L2 name only quarantine. The cumulative
cap is checked as a safety property (it is never reset and runs down only
while detached) and, under fairness, as a liveness property (a detached
session whose cap is used up leaves $\textsc{Det}$, ends or is quarantined);
the model cannot exhibit an unbounded detach/re-attach cycle, so the
absence of one is argued by the cap and shown by the reference bridge and
vectors. Delivery of a drained call across a crash is modelled by the
forwarding acknowledgement: a drained call that was not delivered is
forwarded again after a restart, and a delivered one never is.

\subsection{Attacker model}\label{app:attacker}

This subsection is the threat model of the formal specification; the
broader STRIDE-lite table is in Section~\ref{sec:threat-model}. The model has an
explicit attacker for the re-attach and join path (I11), and TLC checks the
safety properties against it.

\paragraph{Capabilities.} The attacker (i) observes the log and the wire, so
every signature an honest client has produced is known to it; (ii) replays a
request or a signature it has observed, at any later time and in any session
state; (iii) reorders, by presenting a signature after a newer challenge has
been issued or after the nonce it names has been spent or has expired;
(iv) submits requests of its own, with a key that is registered but whose
signing key it does not hold, with a key that is not registered, and with any
nonce number and prefix verdict; and (v) holds signing keys of its own
(\texttt{AdvKeys}), none of which is a client key. It never obtains a signing
key of a registered client; the theft of such a key is out of scope (A7), and
I11-REVOKE, not this model, is the answer to it. It cannot make the bridge
diverge from the log (A4).

\paragraph{Possession.} A signature is a record of the key that signed, the
nonce it covers and the prefix verdict it covers. The predicate
\texttt{Verifies(key, sig, claimed, pre)} holds iff the signing key is the
key the request names and the signed nonce and verdict are the ones the
request carries. An honest client builds only signatures with its own key
over what it presents, and each goes on the wire; the attacker builds
signatures only with its own keys, or resubmits one from the wire. A request
is authentic iff its key is registered and \texttt{Verifies} holds (and, for
I11-REVOKE, the key is not revoked).

\paragraph{Unforgeability assumption (UF).} The model assumes that a
signature verifies for a key only if the holder of that key produced it,
and that it covers the fields \texttt{Verifies} compares. This is the
existential unforgeability of A3 for the signature scheme, stated once and
not model-checked: TLC does not check that Ed25519 or ECDSA P-256 are
unforgeable, and it does not check that the reference bridges sign the
fields the wire format lists (the tests, vectors and mutants do). What TLC
checks is that, if UF holds, the attacker cannot build a signature of its
own that is accepted. The step property
\texttt{NoForgery} (the conjunct on entering \textsc{Res} or joining) states
that every accepted request has a live nonce that it spends and a signature
that an honest client produced for that key, nonce and verdict. The step
properties \texttt{SP\_Honest} and \texttt{SP\_NoReplay} add that an honest
client presented every accepted request, and that a request that was presented
and refused is never accepted later, whatever the registry has become. The
nonce of an honest request for a key that is not yet registered is bound to
that key as not eligible at issue, so adding the key afterwards does not make
an observed signature acceptable.

\paragraph{What TLC checks.} The model records who built a
request. The checks are: a signature replayed under another nonce, or
with another prefix verdict, is refused (a nonce is single use and bound to
the signature); a replayed signature under a spent or expired nonce is
refused and changes nothing (under a live nonce it can be accepted, above); a request signed by a key the attacker holds is
never authentic for a registered key; and the nonce book and the record of
unauthenticated failures stay within their bounds under the same attacker
(I11 (f), (g); the bounds stand for the limits of the reference bridges, which are
configuration values, not protocol constants). Every configuration sets
\texttt{LiveLimit} to one, so the property that a challenge never evicts a
live nonce reads over an empty set and is not meaningfully checked; the
bound on the book is. The model
mutants of the attacker are in the list of mutants above; each is killed.

\paragraph{Limits.} One honest client signs for each key; the attacker
does not interleave sessions, so a signature replayed from one session in
another is not modelled (the model is single-session in its attacker and the
wire format binds the session id). The attacker does not attack the chain
(I5 stays a corruption fault), the transport (I8), the operator signatures
or the attach proof (I-GRANT). Requests whose signature verifies are
recorded without bound, so replaying one observed request fills the log at
the rate of the attacker; I11 (g) bounds only the failures before
authentication, as its text says. The attacker is on only in the attacker
configuration, which has one permission, one call key and no log tampering,
and remembers at most two honest signatures; its checks do not reach the
tool queue, draining, quarantine or an attack that combines three or more
distinct remembered signatures, and the other configurations, including the
two-session one, run without it.

\subsection{Compositional safety arguments}
\label{app:comp-safety}

\paragraph{BC $\circ$ PACE.}
The compositional safety property for PACE is stated over the committed
log of a session, $L$, which is the only store the contract needs:
\[
\forall i:\ L[i].\textit{kind} = \texttt{RESUME\_STARTED}
\Rightarrow
\exists j < i:\ L[j].\textit{kind} = \texttt{RESUME\_PREPARED}
\,\wedge\, L[j].\textit{seq} = L[i].\textit{seq}.
\]
This holds only when the client names the chain head as it stood at
detachment, which is the \texttt{RESUME\_PREPARED} event, or an anchor
that PACE itself defines. It is not claimed for a client that names any
other event of the chain as last seen: the contract admits such a client
(I4) and replays every later event, so \texttt{RESUME\_STARTED}
then names an event that need not be a \texttt{RESUME\_PREPARED}. For the
restricted case, every accepted re-attach names a
\texttt{RESUME\_PREPARED} event that is in the log. Under the contract's own re-attach, a
client is admitted iff its presented history is a prefix of the
hash-chained log and it answers a live challenge, which is guaranteed by
I1 + I4 + I5 + I11 (Lemma~\ref{lem:recover}). This is a statement about
the contract's re-attach, not about PACE: PACE, as read, reserves
\texttt{SESSION\_SUSPEND} (0x05) and \texttt{SESSION\_RESUME} (0x06) in
its message-type registry without a defined body, and defines no resume
credential.

\paragraph{BC $\circ$ GEC.}
The compositional safety property for GEC is vacuous in
GEC-present deployments (the GEC signs events; the bridge forwards
unsigned events). In GEC-less deployments, the bridge is the sole
emitter; the property states that every event in the log was emitted
by the bridge: every event $L[i]$ in the log carries
$L[i].\textit{prev\_hash}=\textit{event\_hash}(L[i-1])$ (I5), so an
event that the bridge did not emit breaks the chain. I-SIGN-EVENT (SHOULD)
governs whether the bridge signs events with an operator key; this is
a deployment choice, not a structural property.

\paragraph{Multi-hop composition (BC $\circ$ BC).}
Multi-hop composition is out of scope (Section~\ref{sec:scope}) and
not modelled here. A future revision may extend the contract with a
\texttt{BC\_HOP} invariant that requires the downstream bridge to
verify the upstream bridge's signature chain before accepting
resumes.

\subsection{Merkle tree for $O(\log n)$ verification (informative)}
\label{app:merkle}

\textbf{Partly normative.} The Merkle tree below is the construction
referenced by the SHOULD-level checkpoint rule I-CHECKPOINT
(Section~\ref{sec:invariants-checkpoint}). The Merkle-augmented
\emph{resume record} $R^+$ remains an informative sketch. The
resume record is bridge-private (I11): a client never holds it, so $R^+$
describes what the bridge could derive for its own use, not a message. The conformance suite in Section~\ref{sec:conformance} still
uses an $O(n)$ full-chain verification.

The conformance suite's full-chain verification is
$O(n)$ in the number of events emitted since session attach. For
long-running sessions ($n > 10^4$), the bridge SHOULD maintain a
Merkle tree over event hashes to enable $O(\log n)$ inclusion proofs.

\begin{definition}[Merkle tree over event log]\label{def:merkle}
Let $\mathit{log} = [e_0, e_1, \ldots, e_{n-1}]$ be the event log of a
session, and let $h_i = \textit{event\_hash}(e_i)$. The Merkle tree
$T(\mathit{log})$ is the Merkle tree hash of RFC~9162, \S2.1.1: the hash of
a leaf is $\textit{SHA-256}(\texttt{0x00} \,\|\, h_i)$, the hash of an
interior node is $\textit{SHA-256}(\texttt{0x01} \,\|\, \textit{left} \,\|\,
\textit{right})$, a tree of $n>1$ leaves is split at the largest power of two
below $n$ (no last node is duplicated), and the root of an empty log is
$\textit{SHA-256}$ of the empty string. The tree's root
$\textit{root}(\mathit{log})$ is the hash of the topmost node. The prefixes
keep a leaf from standing for an interior node.

An inclusion proof for event $e_i$ is the sequence of sibling hashes
along the path from leaf $i$ to the root; verification requires
$O(\log n)$ SHA-256 computations.
\end{definition}

\begin{definition}[Merkle-augmented resume record]\label{def:merkle-rt}
A Merkle-augmented resume record extends
Definition~\ref{def:rt} with two additional fields:
\[
R^+ = R \cup (\textit{merkle\_root}, \textit{inclusion\_proof})
\]
where $\textit{merkle\_root} = \textit{root}(\mathit{log}_{\textit{seq}})$
is the Merkle root over the event log up to and including the
\texttt{RESUME\_PREPARED} event whose sequence number is
$\mathit{seq}$, and $\textit{inclusion\_proof}$ is the inclusion proof
for that event in the Merkle tree. A bridge that keeps $R^+$ verifies
the prefix a re-attaching client names in $O(\log n)$ time.

\end{definition}

The current conformance suite uses an $O(n)$
full-chain verification.

\subsection{Property-based conformance tests}
\label{app:property-tests}

In addition to the example-based tests of
Section~\ref{sec:conformance}, the conformance suite includes the
property-based tests using the Hypothesis library
(\url{https://hypothesis.readthedocs.io/}):

\begin{longtable}{p{0.08\columnwidth} p{0.22\columnwidth} p{0.58\columnwidth}}
\caption{Property-based tests.}\label{tab:property-tests}\\
\toprule
\textbf{Test} & \textbf{Property} & \textbf{Description} \\
\midrule
T8 & Monotonic seq across any action sequence & For any sequence of
bridge operations, the bridge's event log has strictly increasing
\texttt{sequence} numbers (no gaps, no duplicates, no out-of-order). \\
T9 & Invariants across any action sequence & For any sequence
of bridge operations, including restarts and clock jumps, the invariants
I1--I3, I5, I6, I10, I11(d), I14 and I15 hold and the cell equals the fold of the committed
log, at every intermediate state. \\
T10 & Resume from any detached cell & For any sequence that leaves the
session in a cell with $a=\textsc{Det}$, a fresh challenge resumes it and
the invariants hold afterwards. \\
\bottomrule
\end{longtable}

Hypothesis generates random action sequences (up to 40 actions per
sequence for T8 and T9, 30 for T10; 100 examples per property). Each property invariant is checked
after every action, not just at the end. This catches corner cases
that example-based tests miss (e.g., multiple sequential re-attaches
with stale seqs, interleaved detachments and escalations).

\subsection{Traceability matrix}
\label{app:traceability}

Table~\ref{tab:traceability} puts the checks of each invariant and model
rule in one place: the TLA+ property, the tests, the vectors, and the
mutants killed on the Python bridge, on the Rust crate and in the TLA+
module. It is generated by \texttt{supplementary/traceability.py} from the
files that hold those facts (the tags in the pytest docstrings, the
\texttt{covers} field of the vector files, the mutant registries of the Python bridge, the Rust crate and the TLA+ module, with
the verdicts in \texttt{docs/logs/tlc-rerun-2026-11-25/mutants.log}, and the
\texttt{Inv}, \texttt{StepProps} and liveness definitions of
\texttt{BridgeContract.tla}); \texttt{make traceability-check}, part of
\texttt{make validate}, fails on a stale file, on a name that no longer
exists and on an empty cell without a stated limit. An entry
\texttt{Inv.X} is a state invariant and \texttt{StepProps[\ldots]} a
labelled conjunct of the step property, both checked by every safety
configuration of Table~\ref{tab:tlc}; \texttt{L1}, \texttt{L1Res}, \texttt{L2} and
\texttt{CapEnds} are checked by the fair configuration. A ranged entry lists
only members that carry the row. Where a cell says ``---'' it states the
limit that leaves it empty; how far the model covers a row (abstracted flag
or tests only) is argued in Table~\ref{tab:invariant-independence} and not
repeated here. The matrix does not say which configuration reaches which
cell or action; that is Table~\ref{tab:tlc-reach}. Python mutants are filed
under the label their registry gives them, so a name with another prefix
(for example a timer mutant filed under I1) sits in the row of its label's
check. The Rust tests carry no invariant tag, so the matrix lists the Rust
crate through the vectors it runs (all but those that require
\texttt{CHECKPOINT}) and its mutants.

% GENERATED by supplementary/traceability.py --emit-tex; do not edit.
\begingroup\tiny\setlength{\tabcolsep}{2pt}
\begin{longtable}{@{}p{0.09\columnwidth} >{\raggedright\arraybackslash}p{0.17\columnwidth} >{\raggedright\arraybackslash}p{0.14\columnwidth} >{\raggedright\arraybackslash}p{0.12\columnwidth} >{\raggedright\arraybackslash}p{0.44\columnwidth}@{}}
\caption{Traceability matrix: invariant, TLA+ property, tests, vectors and mutants.}\label{tab:traceability}\\
\toprule
\textbf{Row} & \textbf{TLA+ property} & \textbf{Tests} & \textbf{Vectors} & \textbf{Mutants killed (Py: Python bridge; Rs: Rust crate; TLA+: model)} \\
\midrule
\endfirsthead
\toprule
\textbf{Row} & \textbf{TLA+ property} & \textbf{Tests} & \textbf{Vectors} & \textbf{Mutants killed (Py: Python bridge; Rs: Rust crate; TLA+: model)} \\
\midrule
\endhead
I1 & \texttt{Inv.\allowbreak{}I1} & T1, T9, T21b, T45b, T48 & V01, V21b, V23, V23b, V44, V51 & \textbf{Py}: \texttt{no-\allowbreak{}resume-\allowbreak{}prepared}, \texttt{resume-\allowbreak{}record-\allowbreak{}carries-\allowbreak{}a-\allowbreak{}stale-\allowbreak{}scope-\allowbreak{}version}\newline \textbf{Rs}: --- (not among the sampled rules)\newline \textbf{TLA+}: \texttt{I1 detach without a resume record}, \texttt{I1 restart without a resume record} \\
I2 & \texttt{Inv.\allowbreak{}I2}, \texttt{StepProps[I2, I12]} & T2, T9, T25, T25b, T59b, T61b & V12b, V15b, V33, V37, V46 & \textbf{Py}: \texttt{forward-\allowbreak{}while-\allowbreak{}pending}, \texttt{detached-\allowbreak{}call-\allowbreak{}forwarded-\allowbreak{}by-\allowbreak{}default}, \texttt{resolution-\allowbreak{}without-\allowbreak{}an-\allowbreak{}escalation}\newline \textbf{Rs}: --- (not among the sampled rules)\newline \textbf{TLA+}: --- (the model's mutant for this check is filed under I12) \\
I3 & \texttt{StepProps[I3]}, \texttt{StepProps[I3, direct calls]} & T3, T9, T12, T14b--T15, T25b, T56, T56b, T149 & V04, V06--V06c, V15b, V18, V27, V27b, V47 & \textbf{Py}: \texttt{accept-\allowbreak{}widening-\allowbreak{}detached}, \texttt{narrow-\allowbreak{}accepts-\allowbreak{}widening}, \texttt{observed-\allowbreak{}widening-\allowbreak{}recorded}, \texttt{observed-\allowbreak{}widening-\allowbreak{}not-\allowbreak{}reverted}, \texttt{scope-\allowbreak{}recorded-\allowbreak{}as-\allowbreak{}written-\allowbreak{}at-\allowbreak{}attach}, \texttt{direct-\allowbreak{}call-\allowbreak{}outside-\allowbreak{}the-\allowbreak{}scope-\allowbreak{}forwarded}, \texttt{direct-\allowbreak{}refusal-\allowbreak{}not-\allowbreak{}recorded}, \texttt{direct-\allowbreak{}refusal-\allowbreak{}removes-\allowbreak{}a-\allowbreak{}queued-\allowbreak{}call}, \texttt{observed-\allowbreak{}widening-\allowbreak{}refused-\allowbreak{}when-\allowbreak{}key-\allowbreak{}exhausted}\newline \textbf{Rs}: --- (not among the sampled rules)\newline \textbf{TLA+}: \texttt{I3 widening applied}, \texttt{I3 a direct call outside the scope is forwarded}, \texttt{I3 a direct call outside the scope is refused without a record} \\
I4 & \texttt{StepProps[I4, I9, I11, I15]} & T1, T7, T7b, T11, T19, T52 & V01, V02, V21, V44, V45 & \textbf{Py}: \texttt{no-\allowbreak{}prefix-\allowbreak{}check}, \texttt{no-\allowbreak{}replay}, \texttt{replay-\allowbreak{}off-\allowbreak{}by-\allowbreak{}one}, \texttt{negative-\allowbreak{}sequence-\allowbreak{}number-\allowbreak{}accepted}\newline \textbf{Rs}: \texttt{PREFIX-\allowbreak{}hash-\allowbreak{}unchecked}\newline \textbf{TLA+}: \texttt{I4 the prefix is not checked} \\
I5 & \texttt{StepProps[I5]} & T1, T5, T5b, T9, T22, T38, T66d, T66e & V01, V19, V29, V44, V54, V54c & \textbf{Py}: \texttt{no-\allowbreak{}prev-\allowbreak{}hash}, \texttt{resume-\allowbreak{}failure-\allowbreak{}terminates}, \texttt{recovery-\allowbreak{}failure-\allowbreak{}terminates}, \texttt{quarantined-\allowbreak{}session-\allowbreak{}still-\allowbreak{}served}, \texttt{quarantine-\allowbreak{}not-\allowbreak{}durable}, \texttt{chain-\allowbreak{}check-\allowbreak{}ignores-\allowbreak{}the-\allowbreak{}sequence-\allowbreak{}number}\newline \textbf{Rs}: \texttt{CHAIN-\allowbreak{}writer-\allowbreak{}drops-\allowbreak{}prev-\allowbreak{}hash}, \texttt{CHAIN-\allowbreak{}verifier-\allowbreak{}ignores-\allowbreak{}prev-\allowbreak{}hash}, \texttt{CHAIN-\allowbreak{}verifier-\allowbreak{}ignores-\allowbreak{}sequence}, \texttt{CANON-\allowbreak{}non-\allowbreak{}ascii-\allowbreak{}escaped}, \texttt{CANON-\allowbreak{}keys-\allowbreak{}by-\allowbreak{}code-\allowbreak{}point}, \texttt{CANON-\allowbreak{}float-\allowbreak{}printed-\allowbreak{}natively}, \texttt{CANON-\allowbreak{}large-\allowbreak{}integer-\allowbreak{}accepted}\newline \textbf{TLA+}: \texttt{I5 resume verifies a corrupt chain} \\
I6 & \texttt{Inv.\allowbreak{}I6}, \texttt{StepProps[I6]} & T4, T6, T9, T40, T66b, T66e, T145 & V49, V50, V54b, V118, V130 & \textbf{Py}: \texttt{ops-\allowbreak{}after-\allowbreak{}termination}\newline \textbf{Rs}: --- (not among the sampled rules)\newline \textbf{TLA+}: \texttt{I6 a termination keeps the queue}, \texttt{I6 a termination leaves the timers set}, \texttt{I6 an escalation is accepted after termination}, \texttt{I6 time runs on after termination} \\
I7 & \texttt{StepProps[I7]} & T1, T6, T11, T32, T39 & V01, V03 & \textbf{Py}: \texttt{mutate-\allowbreak{}before-\allowbreak{}append}\newline \textbf{Rs}: --- (not among the sampled rules)\newline \textbf{TLA+}: \texttt{I7 cell change without an event} \\
I8 & --- (transport flag; not modelled) & T16, T152g & V09, V09b & \textbf{Py}: \texttt{ignore-\allowbreak{}replay}, \texttt{flagged-\allowbreak{}request-\allowbreak{}keeps-\allowbreak{}its-\allowbreak{}nonce}\newline \textbf{Rs}: --- (not among the sampled rules)\newline \textbf{TLA+}: --- (not modelled) \\
I9 & \texttt{Inv.\allowbreak{}I9}, \texttt{StepProps[I4, I9, I11, I15]}, \texttt{StepProps[I9/I11]} & T13, T13b, T31 & V05, V42 & \textbf{Py}: \texttt{ignore-\allowbreak{}expiry}, \texttt{expiry-\allowbreak{}terminates}\newline \textbf{Rs}: \texttt{NONCE-\allowbreak{}expiry-\allowbreak{}ignored}\newline \textbf{TLA+}: \texttt{I9 an expired nonce is accepted}, \texttt{I9 a nonce outlives T\_\allowbreak{}nonce} \\
I10 & \texttt{Inv.\allowbreak{}I10}, \texttt{StepProps[I10]} & T9, T14--T14c, T15b, T45b, T50 & V06--V06c, V15, V21, V21b, V47 & \textbf{Py}: \texttt{version-\allowbreak{}not-\allowbreak{}incremented}, \texttt{noop-\allowbreak{}narrowing-\allowbreak{}accepted}, \texttt{narrowing-\allowbreak{}refusal-\allowbreak{}not-\allowbreak{}named}, \texttt{narrowed-\allowbreak{}scope-\allowbreak{}recorded-\allowbreak{}as-\allowbreak{}written}\newline \textbf{Rs}: --- (not among the sampled rules)\newline \textbf{TLA+}: \texttt{I10 narrowing without version bump}, \texttt{I10 the version advances without a narrowing} \\
I11 (a)--(e) & \texttt{Inv.\allowbreak{}I11}, \texttt{StepProps[I4, I9, I11, I15]}, \texttt{StepProps[I11, honest presenter]}, \texttt{StepProps[I11, no replay]}, \texttt{StepProps[I11d]}, \texttt{StepProps[I9/I11]} & T1, T7, T9, T13b, T17--T19b, T19d, T26b, T35, T51, T69b, T71b, T82, T85, T103--T112, T118, T121, T126, T127, T141--T144, T147, T150b, T151, T152--T152g, T156, T157 & V01, V02, V07, V08, V09b, V16b--V16d, V42, V44, V45, V59--V62f, V70, V86, V87, V90, V90b, V90d--V98, V101--V110, V114, V127--V129, V131, V133, V134 & \textbf{Py}: \texttt{drop-\allowbreak{}proof-\allowbreak{}check}, \texttt{resume-\allowbreak{}domain-\allowbreak{}tag-\allowbreak{}dropped}, \texttt{resume-\allowbreak{}key-\allowbreak{}id-\allowbreak{}dropped}, \texttt{resume-\allowbreak{}operation-\allowbreak{}ignored}, \texttt{join-\allowbreak{}tagged-\allowbreak{}as-\allowbreak{}reattach}, \texttt{scheme-\allowbreak{}tag-\allowbreak{}ignored}, \texttt{nonce-\allowbreak{}not-\allowbreak{}single-\allowbreak{}use}, \texttt{nonce-\allowbreak{}spent-\allowbreak{}only-\allowbreak{}on-\allowbreak{}success}, \texttt{unauthentic-\allowbreak{}request-\allowbreak{}spends-\allowbreak{}nonce}, \texttt{unknown-\allowbreak{}key-\allowbreak{}spends-\allowbreak{}the-\allowbreak{}nonce}, \texttt{nonce-\allowbreak{}not-\allowbreak{}bound-\allowbreak{}to-\allowbreak{}the-\allowbreak{}key-\allowbreak{}id}, \texttt{nonce-\allowbreak{}issued-\allowbreak{}before-\allowbreak{}registration-\allowbreak{}verifies}, \texttt{eligibility-\allowbreak{}read-\allowbreak{}when-\allowbreak{}presented}, \texttt{binding-\allowbreak{}not-\allowbreak{}checked}, \texttt{binding-\allowbreak{}refusal-\allowbreak{}spends-\allowbreak{}the-\allowbreak{}nonce}, \texttt{binding-\allowbreak{}refusal-\allowbreak{}distinguishable}, \texttt{challenge-\allowbreak{}refuses-\allowbreak{}an-\allowbreak{}unregistered-\allowbreak{}key-\allowbreak{}id}, \texttt{registry-\allowbreak{}stores-\allowbreak{}signing-\allowbreak{}seed}, \texttt{added-\allowbreak{}key-\allowbreak{}registry-\allowbreak{}stores-\allowbreak{}signing-\allowbreak{}seed}, \texttt{nonce-\allowbreak{}not-\allowbreak{}bound-\allowbreak{}to-\allowbreak{}session}, \texttt{bad-\allowbreak{}reattach-\allowbreak{}terminates}, \texttt{registry-\allowbreak{}global-\allowbreak{}by-\allowbreak{}key-\allowbreak{}id}, \texttt{registry-\allowbreak{}reregistration-\allowbreak{}allowed}, \texttt{registry-\allowbreak{}recovery-\allowbreak{}leaves-\allowbreak{}differing-\allowbreak{}entries}, \texttt{registry-\allowbreak{}recovery-\allowbreak{}leaves-\allowbreak{}extra-\allowbreak{}entries}, \texttt{registry-\allowbreak{}recovery-\allowbreak{}leaves-\allowbreak{}missing-\allowbreak{}entries}, \texttt{registry-\allowbreak{}recovery-\allowbreak{}repair-\allowbreak{}unreported}, \texttt{registry-\allowbreak{}recovery-\allowbreak{}repairs-\allowbreak{}a-\allowbreak{}broken-\allowbreak{}chain-\allowbreak{}session}, \texttt{registry-\allowbreak{}client-\allowbreak{}added-\allowbreak{}logs-\allowbreak{}wrong-\allowbreak{}key}, \texttt{registry-\allowbreak{}attached-\allowbreak{}logs-\allowbreak{}no-\allowbreak{}scheme}, \texttt{registry-\allowbreak{}failed-\allowbreak{}attach-\allowbreak{}keeps-\allowbreak{}entry}, \texttt{session-\allowbreak{}nonce-\allowbreak{}book-\allowbreak{}shared}, \texttt{reattach-\allowbreak{}on-\allowbreak{}attached-\allowbreak{}session-\allowbreak{}joins}, \texttt{already-\allowbreak{}attached-\allowbreak{}before-\allowbreak{}authentication}, \texttt{join-\allowbreak{}attachment-\allowbreak{}answered-\allowbreak{}before-\allowbreak{}authentication}, \texttt{join-\allowbreak{}refused-\allowbreak{}as-\allowbreak{}already-\allowbreak{}attached}, \texttt{attach-\allowbreak{}accepts-\allowbreak{}signing-\allowbreak{}key}, \texttt{challenge-\allowbreak{}leaks-\allowbreak{}state}, \texttt{reattach-\allowbreak{}leaks-\allowbreak{}state}, \texttt{add-\allowbreak{}client-\allowbreak{}accepts-\allowbreak{}signing-\allowbreak{}key}, \texttt{unknown-\allowbreak{}key-\allowbreak{}distinguishable}\newline \textbf{Rs}: \texttt{SIG-\allowbreak{}skip-\allowbreak{}check}, \texttt{SIG-\allowbreak{}ed25519-\allowbreak{}accepts-\allowbreak{}any}, \texttt{SIG-\allowbreak{}p256-\allowbreak{}accepts-\allowbreak{}any}, \texttt{CHALLENGE-\allowbreak{}leaks-\allowbreak{}state}, \texttt{RESUME-\allowbreak{}leaks-\allowbreak{}state}, \texttt{NONCE-\allowbreak{}not-\allowbreak{}spent}, \texttt{REVOKE-\allowbreak{}nonce-\allowbreak{}not-\allowbreak{}spent}, \texttt{RESUME-\allowbreak{}domain-\allowbreak{}tag-\allowbreak{}dropped}, \texttt{RESUME-\allowbreak{}key-\allowbreak{}id-\allowbreak{}dropped}, \texttt{RESUME-\allowbreak{}join-\allowbreak{}verified-\allowbreak{}as-\allowbreak{}reattach}, \texttt{RESUME-\allowbreak{}join-\allowbreak{}tagged-\allowbreak{}as-\allowbreak{}reattach}, \texttt{REGISTRY-\allowbreak{}global-\allowbreak{}by-\allowbreak{}key-\allowbreak{}id}, \texttt{REGISTRY-\allowbreak{}reregistration-\allowbreak{}allowed}, \texttt{RECOVERY-\allowbreak{}malformed-\allowbreak{}log-\allowbreak{}not-\allowbreak{}isolated}, \texttt{REVOKE-\allowbreak{}requests-\allowbreak{}unbounded}, \texttt{REGISTRY-\allowbreak{}recovery-\allowbreak{}leaves-\allowbreak{}differing-\allowbreak{}entries}, \texttt{REGISTRY-\allowbreak{}recovery-\allowbreak{}leaves-\allowbreak{}extra-\allowbreak{}entries}, \texttt{REGISTRY-\allowbreak{}recovery-\allowbreak{}leaves-\allowbreak{}missing-\allowbreak{}entries}, \texttt{REGISTRY-\allowbreak{}recovery-\allowbreak{}repair-\allowbreak{}unreported}, \texttt{REGISTRY-\allowbreak{}client-\allowbreak{}added-\allowbreak{}logs-\allowbreak{}wrong-\allowbreak{}key}, \texttt{REGISTRY-\allowbreak{}failed-\allowbreak{}attach-\allowbreak{}keeps-\allowbreak{}entry}, \texttt{RESUME-\allowbreak{}reattach-\allowbreak{}on-\allowbreak{}attached-\allowbreak{}session-\allowbreak{}joins}, \texttt{RESUME-\allowbreak{}join-\allowbreak{}refused-\allowbreak{}as-\allowbreak{}already-\allowbreak{}attached}\newline \textbf{TLA+}: \texttt{I11 an unauthentic request is accepted}, \texttt{I11 a spent nonce is accepted again}, \texttt{I11 an accepted nonce is not consumed}, \texttt{I11 a nonce issued for an unregistered key is marked good}, \texttt{I11 a challenge is refused for an unregistered key}, \texttt{I11 a nonce issued before its key was registered verifies afterwards}, \texttt{I11 a replayed request is accepted once its key is registered}, \texttt{I11 a nonce is accepted for a key it was not issued for}, \texttt{I11 CLIENT\_\allowbreak{}ADDED revives the nonces issued for the new key}, \texttt{I11 OPERATOR\_\allowbreak{}ADD\_\allowbreak{}KEY revives the nonces issued for the new key}, \texttt{I11 a request refused during a verification keeps its nonce}, \texttt{I11 a restart forgets which nonces were spent}, \texttt{I11 a detached session is attached without a request}, \texttt{Attacker holds a client's signing key}, \texttt{Attacker: a signature verifies under any key}, \texttt{Attacker: a signature is not bound to its nonce} \\
I11 (f) & \texttt{Inv.\allowbreak{}IBook}, \texttt{StepProps[I11f]} & T97, T98, T100, T152e, T152h & V62e, V86 & \textbf{Py}: \texttt{I11-\allowbreak{}challenge-\allowbreak{}for-\allowbreak{}unregistered-\allowbreak{}key-\allowbreak{}id-\allowbreak{}is-\allowbreak{}free}, \texttt{I11-\allowbreak{}challenge-\allowbreak{}evicts-\allowbreak{}live-\allowbreak{}nonce}\newline \textbf{Rs}: \texttt{NONCE-\allowbreak{}challenge-\allowbreak{}evicts-\allowbreak{}live-\allowbreak{}nonce}, \texttt{SESSION-\allowbreak{}challenge-\allowbreak{}budget-\allowbreak{}shared}\newline \textbf{TLA+}: \texttt{I11f a challenge ignores the live-\allowbreak{}nonce limit}, \texttt{I11f a challenge evicts a live nonce}, \texttt{I11f a challenge for an unregistered key issues a dead nonce} \\
I11 (g) & \texttt{Inv.\allowbreak{}IFail}, \texttt{StepProps[I11g]} & T99--T100, T142, T152h, T153 & V70b, V87 & \textbf{Py}: \texttt{session-\allowbreak{}unauth-\allowbreak{}budget-\allowbreak{}shared}, \texttt{session-\allowbreak{}unauth-\allowbreak{}budget-\allowbreak{}restored-\allowbreak{}for-\allowbreak{}all}, \texttt{I11-\allowbreak{}unauthenticated-\allowbreak{}rejections-\allowbreak{}unbounded}, \texttt{I11-\allowbreak{}revoked-\allowbreak{}key-\allowbreak{}refusals-\allowbreak{}logged-\allowbreak{}without-\allowbreak{}bound}, \texttt{I11-\allowbreak{}binding-\allowbreak{}refusals-\allowbreak{}logged-\allowbreak{}without-\allowbreak{}bound}\newline \textbf{Rs}: \texttt{LOG-\allowbreak{}unauthenticated-\allowbreak{}rejections-\allowbreak{}unbounded}, \texttt{SESSION-\allowbreak{}unauth-\allowbreak{}budget-\allowbreak{}restored-\allowbreak{}for-\allowbreak{}all}\newline \textbf{TLA+}: \texttt{I11g unauthenticated failures are all recorded}, \texttt{I11g a request that verifies is dropped past the limit}, \texttt{I11g an authentic request does not restore the allowance} \\
I12 & \texttt{StepProps[I2, I12]} & T32--T34, T38, T39, T88, T136, T143 & V10, V11, V19, V32, V36, V38, V39, V40, V43, V48, V71, V72, V102--V104, V113, V116, V129 & \textbf{Py}: \texttt{forward-\allowbreak{}before-\allowbreak{}commit}, \texttt{recovery-\allowbreak{}skips-\allowbreak{}replay}\newline \textbf{Rs}: \texttt{EXHAUST-\allowbreak{}restart-\allowbreak{}skipped}\newline \textbf{TLA+}: \texttt{I12 forward before the DRAINED event} \\
I13 & \texttt{Inv.\allowbreak{}I13}, \texttt{StepProps[Rule D, I13]}, \texttt{StepProps[I13, approval]} & T2, T41--T45, T60e, T61b, T154 & V12--V15, V35, V38, V38b, V41--V41d, V46, V48 & \textbf{Py}: \texttt{no-\allowbreak{}queue-\allowbreak{}limit}, \texttt{duplicate-\allowbreak{}keys-\allowbreak{}queued}, \texttt{no-\allowbreak{}approval-\allowbreak{}expiry}, \texttt{approval-\allowbreak{}valid-\allowbreak{}at-\allowbreak{}t-\allowbreak{}approve}, \texttt{later-\allowbreak{}approval-\allowbreak{}revives-\allowbreak{}a-\allowbreak{}lapsed-\allowbreak{}one}, \texttt{no-\allowbreak{}scope-\allowbreak{}revalidation}, \texttt{no-\allowbreak{}call-\allowbreak{}expiry}, \texttt{acknowledged-\allowbreak{}call-\allowbreak{}forwarded-\allowbreak{}again}, \texttt{generated-\allowbreak{}key-\allowbreak{}collides-\allowbreak{}with-\allowbreak{}a-\allowbreak{}supplied-\allowbreak{}one}, \texttt{second-\allowbreak{}completion-\allowbreak{}accepted}\newline \textbf{Rs}: --- (not among the sampled rules)\newline \textbf{TLA+}: \texttt{I13 drain skips scope revalidation}, \texttt{I13 drain ignores approval expiry}, \texttt{I13 a later resolution renews a lapsed approval}, \texttt{I13 queue bound ignored} \\
I14 & \texttt{StepProps[Rule D, I13]} & T9, T53--T55, T58, T129, T130, T158 & V12, V25, V26, V34, V41, V41b, V46, V112, V113 & \textbf{Py}: \texttt{hash-\allowbreak{}equality-\allowbreak{}not-\allowbreak{}checked}, \texttt{unanswered-\allowbreak{}call-\allowbreak{}drained}, \texttt{denial-\allowbreak{}ignored}, \texttt{second-\allowbreak{}answer-\allowbreak{}accepted}, \texttt{shown-\allowbreak{}hash-\allowbreak{}not-\allowbreak{}of-\allowbreak{}content}, \texttt{drained-\allowbreak{}omits-\allowbreak{}hash}, \texttt{hash-\allowbreak{}omits-\allowbreak{}tool-\allowbreak{}and-\allowbreak{}permission}, \texttt{forwards-\allowbreak{}without-\allowbreak{}hashed-\allowbreak{}content}, \texttt{queued-\allowbreak{}content-\allowbreak{}not-\allowbreak{}retained}, \texttt{held-\allowbreak{}call-\allowbreak{}forwarded-\allowbreak{}unanswered}\newline \textbf{Rs}: \texttt{CALL-\allowbreak{}hash-\allowbreak{}omits-\allowbreak{}tool-\allowbreak{}and-\allowbreak{}permission}, \texttt{CALL-\allowbreak{}forwards-\allowbreak{}without-\allowbreak{}hashed-\allowbreak{}content}\newline \textbf{TLA+}: \texttt{I14 drain ignores the answer} \\
I15 & \texttt{Inv.\allowbreak{}I15}, \texttt{StepProps[I4, I9, I11, I15]}, \texttt{StepProps[I15d]}, \texttt{StepProps[I15b]} & T9, T18b, T19d, T21b, T47--T50, T86, T119--T123, T141, T141b, T147, T152, T152b, T155 & V22--V24, V45, V51, V62, V62b, V62d, V62f, V63, V65, V74, V90--V90d, V95, V96, V105--V108b, V111c, V120, V127, V128, V131, V133, V134 & \textbf{Py}: \texttt{partial-\allowbreak{}departure-\allowbreak{}detaches}, \texttt{unregistered-\allowbreak{}key-\allowbreak{}accepted}, \texttt{registration-\allowbreak{}needs-\allowbreak{}no-\allowbreak{}client}, \texttt{add-\allowbreak{}client-\allowbreak{}skip-\allowbreak{}verify}, \texttt{add-\allowbreak{}client-\allowbreak{}nonce-\allowbreak{}not-\allowbreak{}spent}, \texttt{add-\allowbreak{}client-\allowbreak{}nonce-\allowbreak{}not-\allowbreak{}bound-\allowbreak{}to-\allowbreak{}the-\allowbreak{}key-\allowbreak{}id}, \texttt{add-\allowbreak{}client-\allowbreak{}connection-\allowbreak{}answered-\allowbreak{}before-\allowbreak{}the-\allowbreak{}signature}, \texttt{add-\allowbreak{}client-\allowbreak{}domain-\allowbreak{}tag-\allowbreak{}dropped}, \texttt{add-\allowbreak{}client-\allowbreak{}public-\allowbreak{}key-\allowbreak{}unbound}, \texttt{add-\allowbreak{}client-\allowbreak{}revoked-\allowbreak{}key-\allowbreak{}authorises}, \texttt{join-\allowbreak{}does-\allowbreak{}not-\allowbreak{}add-\allowbreak{}client}, \texttt{resume-\allowbreak{}attaches-\allowbreak{}nobody}, \texttt{added-\allowbreak{}key-\allowbreak{}not-\allowbreak{}registered}, \texttt{join-\allowbreak{}proof-\allowbreak{}starts-\allowbreak{}a-\allowbreak{}reattach}, \texttt{foreign-\allowbreak{}disconnect-\allowbreak{}abandons-\allowbreak{}the-\allowbreak{}reattach}\newline \textbf{Rs}: \texttt{KEY-\allowbreak{}unregistered-\allowbreak{}accepted}, \texttt{ADDCLIENT-\allowbreak{}skip-\allowbreak{}verify}, \texttt{ADDCLIENT-\allowbreak{}nonce-\allowbreak{}not-\allowbreak{}spent}, \texttt{ADDCLIENT-\allowbreak{}domain-\allowbreak{}tag-\allowbreak{}dropped}\newline \textbf{TLA+}: \texttt{I15 a partial departure detaches}, \texttt{I15 a partial departure leaves the attached cell}, \texttt{I15 a key is registered without a connected client}, \texttt{I15 a completed resume attaches nobody}, \texttt{I15 a restart keeps the client set} \\
I11-REVOKE & \texttt{Inv.\allowbreak{}IRev}, \texttt{StepProps[I11-\allowbreak{}REVOKE]}, \texttt{StepProps[I11-\allowbreak{}REVOKE]}, \texttt{StepProps[key\_\allowbreak{}exhausted]} & T56b, T67--T76, T85--T89b, T99b, T122, T128, T128b, T132, T134--T137, T145--T146b, T150, T151, T152d, T153 & V27b, V62c, V63--V85, V90c, V111--V111c, V114b--V126, V130, V132 & \textbf{Py}: \texttt{session-\allowbreak{}revocation-\allowbreak{}requests-\allowbreak{}shared}, \texttt{session-\allowbreak{}revocation-\allowbreak{}requests-\allowbreak{}restored-\allowbreak{}for-\allowbreak{}all}, \texttt{accept-\allowbreak{}revoked-\allowbreak{}key}, \texttt{revoked-\allowbreak{}request-\allowbreak{}keeps-\allowbreak{}nonce}, \texttt{revoked-\allowbreak{}key-\allowbreak{}spends-\allowbreak{}another-\allowbreak{}keys-\allowbreak{}nonce}, \texttt{soft-\allowbreak{}detaches}, \texttt{res-\allowbreak{}not-\allowbreak{}abandoned}, \texttt{soft-\allowbreak{}revocation-\allowbreak{}leaves-\allowbreak{}the-\allowbreak{}verification}, \texttt{revocation-\allowbreak{}checked-\allowbreak{}before-\allowbreak{}the-\allowbreak{}signature}, \texttt{revoked-\allowbreak{}key-\allowbreak{}may-\allowbreak{}join}, \texttt{free-\allowbreak{}form-\allowbreak{}reason-\allowbreak{}accepted}, \texttt{key-\allowbreak{}revoked-\allowbreak{}twice}, \texttt{soft-\allowbreak{}leaves-\allowbreak{}client-\allowbreak{}set}, \texttt{hard-\allowbreak{}keeps-\allowbreak{}client-\allowbreak{}set}, \texttt{key-\allowbreak{}exhausted-\allowbreak{}not-\allowbreak{}blocked}, \texttt{TIMERS-\allowbreak{}key-\allowbreak{}exhausted-\allowbreak{}timers-\allowbreak{}frozen}, \texttt{key-\allowbreak{}exhausted-\allowbreak{}restart-\allowbreak{}skipped}, \texttt{key-\allowbreak{}exhausted-\allowbreak{}refused-\allowbreak{}before-\allowbreak{}timers}, \texttt{key-\allowbreak{}exhausted-\allowbreak{}disconnect-\allowbreak{}refused}, \texttt{add-\allowbreak{}key-\allowbreak{}after-\allowbreak{}expiry}, \texttt{exhaustion-\allowbreak{}not-\allowbreak{}derived}, \texttt{add-\allowbreak{}key-\allowbreak{}not-\allowbreak{}replayed}, \texttt{add-\allowbreak{}key-\allowbreak{}readmits-\allowbreak{}registered-\allowbreak{}id}, \texttt{add-\allowbreak{}key-\allowbreak{}outside-\allowbreak{}exhaustion}, \texttt{added-\allowbreak{}key-\allowbreak{}not-\allowbreak{}registered}, \texttt{terminate-\allowbreak{}outside-\allowbreak{}exhaustion}, \texttt{unknown-\allowbreak{}key-\allowbreak{}revocable}, \texttt{request-\allowbreak{}executes-\allowbreak{}revocation}, \texttt{absent-\allowbreak{}signature-\allowbreak{}is-\allowbreak{}not-\allowbreak{}a-\allowbreak{}request}, \texttt{requests-\allowbreak{}unbounded}, \texttt{bare-\allowbreak{}terminate-\allowbreak{}ends-\allowbreak{}exhausted-\allowbreak{}session}, \texttt{exhausted-\allowbreak{}refusal-\allowbreak{}blocks-\allowbreak{}quarantine-\allowbreak{}exit}, \texttt{exhausted-\allowbreak{}terminate-\allowbreak{}skips-\allowbreak{}timers}\newline \textbf{Rs}: \texttt{REVOKE-\allowbreak{}revoked-\allowbreak{}key-\allowbreak{}accepted}, \texttt{EXHAUST-\allowbreak{}bare-\allowbreak{}terminate-\allowbreak{}ends-\allowbreak{}session}, \texttt{EXHAUST-\allowbreak{}bare-\allowbreak{}terminate-\allowbreak{}blocks-\allowbreak{}quarantine-\allowbreak{}exit}, \texttt{EXHAUST-\allowbreak{}bare-\allowbreak{}terminate-\allowbreak{}skips-\allowbreak{}timers}, \texttt{SESSION-\allowbreak{}revocation-\allowbreak{}requests-\allowbreak{}shared}, \texttt{SESSION-\allowbreak{}revocation-\allowbreak{}requests-\allowbreak{}restored-\allowbreak{}for-\allowbreak{}all}\newline \textbf{TLA+}: \texttt{I11-\allowbreak{}REVOKE a revoked key is accepted}, \texttt{I11-\allowbreak{}REVOKE revoking a key removes it from K}, \texttt{I11-\allowbreak{}REVOKE a key outside K can be revoked}, \texttt{I11-\allowbreak{}REVOKE a hard revocation leaves the client connected}, \texttt{I11-\allowbreak{}REVOKE revoking the last client does not detach}, \texttt{I11-\allowbreak{}REVOKE a hard revocation does not abandon a verification}, \texttt{I11-\allowbreak{}REVOKE a soft revocation does not abandon a verification}, \texttt{I11-\allowbreak{}REVOKE a revoked client registers a key}, \texttt{key\_\allowbreak{}exhausted does not freeze the session}, \texttt{key\_\allowbreak{}exhausted freezes the timers}, \texttt{key\_\allowbreak{}exhausted: a timeout does not fire} \\
I-OPERATOR & \texttt{StepProps[key\_\allowbreak{}exhausted]} & T90--T96, T138--T140, T150, T150b & V75--V85, V114b, V121--V126, V132, V132b & \textbf{Py}: \texttt{revoke-\allowbreak{}skip-\allowbreak{}verify}, \texttt{revoke-\allowbreak{}names-\allowbreak{}the-\allowbreak{}key-\allowbreak{}before-\allowbreak{}the-\allowbreak{}signature}, \texttt{terminate-\allowbreak{}names-\allowbreak{}the-\allowbreak{}session-\allowbreak{}before-\allowbreak{}the-\allowbreak{}signature}, \texttt{add-\allowbreak{}key-\allowbreak{}names-\allowbreak{}the-\allowbreak{}state-\allowbreak{}before-\allowbreak{}the-\allowbreak{}signature}, \texttt{terminate-\allowbreak{}of-\allowbreak{}a-\allowbreak{}terminated-\allowbreak{}session-\allowbreak{}succeeds}, \texttt{add-\allowbreak{}key-\allowbreak{}skip-\allowbreak{}verify}, \texttt{terminate-\allowbreak{}skip-\allowbreak{}verify}, \texttt{no-\allowbreak{}verifier-\allowbreak{}accepts}, \texttt{missing-\allowbreak{}signature-\allowbreak{}accepted}, \texttt{revoke-\allowbreak{}not-\allowbreak{}bound-\allowbreak{}to-\allowbreak{}reason}, \texttt{revoke-\allowbreak{}not-\allowbreak{}bound-\allowbreak{}to-\allowbreak{}key}, \texttt{add-\allowbreak{}key-\allowbreak{}not-\allowbreak{}bound-\allowbreak{}to-\allowbreak{}key-\allowbreak{}material}, \texttt{terminate-\allowbreak{}not-\allowbreak{}bound-\allowbreak{}to-\allowbreak{}session}, \texttt{revoke-\allowbreak{}not-\allowbreak{}bound-\allowbreak{}to-\allowbreak{}bridge}, \texttt{add-\allowbreak{}key-\allowbreak{}not-\allowbreak{}bound-\allowbreak{}to-\allowbreak{}bridge}, \texttt{terminate-\allowbreak{}not-\allowbreak{}bound-\allowbreak{}to-\allowbreak{}bridge}, \texttt{revoke-\allowbreak{}not-\allowbreak{}bound-\allowbreak{}to-\allowbreak{}position}, \texttt{add-\allowbreak{}key-\allowbreak{}not-\allowbreak{}bound-\allowbreak{}to-\allowbreak{}position}, \texttt{terminate-\allowbreak{}not-\allowbreak{}bound-\allowbreak{}to-\allowbreak{}position}, \texttt{revoke-\allowbreak{}binds-\allowbreak{}sequence-\allowbreak{}only}, \texttt{revoke-\allowbreak{}binds-\allowbreak{}hash-\allowbreak{}only}, \texttt{domain-\allowbreak{}tags-\allowbreak{}collide}\newline \textbf{Rs}: \texttt{OPERATOR-\allowbreak{}revoke-\allowbreak{}omits-\allowbreak{}bridge-\allowbreak{}id}, \texttt{OPERATOR-\allowbreak{}add-\allowbreak{}key-\allowbreak{}omits-\allowbreak{}bridge-\allowbreak{}id}, \texttt{OPERATOR-\allowbreak{}terminate-\allowbreak{}omits-\allowbreak{}bridge-\allowbreak{}id}, \texttt{OPERATOR-\allowbreak{}revoke-\allowbreak{}omits-\allowbreak{}position}, \texttt{OPERATOR-\allowbreak{}add-\allowbreak{}key-\allowbreak{}omits-\allowbreak{}position}, \texttt{OPERATOR-\allowbreak{}terminate-\allowbreak{}omits-\allowbreak{}position}, \texttt{OPERATOR-\allowbreak{}revoke-\allowbreak{}binds-\allowbreak{}sequence-\allowbreak{}only}\newline \textbf{TLA+}: \texttt{OPERATOR\_\allowbreak{}ADD\_\allowbreak{}KEY on a session that is not key\_\allowbreak{}exhausted}, \texttt{OPERATOR\_\allowbreak{}ADD\_\allowbreak{}KEY clears the revocation markers} \\
I-GRANT & --- (attach and grants not modelled) & T77--T81, T113--T118, T124 & V55--V58, V97--V101 & \textbf{Py}: \texttt{skip-\allowbreak{}verify}, \texttt{no-\allowbreak{}verifier-\allowbreak{}accepts}, \texttt{grant-\allowbreak{}not-\allowbreak{}bound-\allowbreak{}to-\allowbreak{}scope}, \texttt{grant-\allowbreak{}omits-\allowbreak{}public-\allowbreak{}key}, \texttt{grant-\allowbreak{}omits-\allowbreak{}bridge-\allowbreak{}id}, \texttt{attach-\allowbreak{}skip-\allowbreak{}possession}, \texttt{possession-\allowbreak{}omits-\allowbreak{}bridge-\allowbreak{}id}, \texttt{possession-\allowbreak{}uses-\allowbreak{}grant-\allowbreak{}domain}, \texttt{possession-\allowbreak{}omits-\allowbreak{}public-\allowbreak{}key}\newline \textbf{Rs}: \texttt{ATTACH-\allowbreak{}possession-\allowbreak{}not-\allowbreak{}checked}, \texttt{GRANT-\allowbreak{}omits-\allowbreak{}public-\allowbreak{}key}, \texttt{GRANT-\allowbreak{}omits-\allowbreak{}bridge-\allowbreak{}id}, \texttt{ATTACH-\allowbreak{}proof-\allowbreak{}omits-\allowbreak{}bridge-\allowbreak{}id}, \texttt{ATTACH-\allowbreak{}proof-\allowbreak{}uses-\allowbreak{}grant-\allowbreak{}domain}\newline \textbf{TLA+}: --- (not modelled) \\
I-CHECKPOINT & --- (not modelled) & T63--T66f & V52--V54c, V88, V89 & \textbf{Py}: \texttt{ICHECK-\allowbreak{}drop-\allowbreak{}checkpoint-\allowbreak{}emission}, \texttt{ICHECK-\allowbreak{}wrong-\allowbreak{}root}, \texttt{ICHECK-\allowbreak{}count-\allowbreak{}trigger-\allowbreak{}total-\allowbreak{}length}, \texttt{ICHECK-\allowbreak{}missed-\allowbreak{}deadline}, \texttt{ICHECK-\allowbreak{}checkpoint-\allowbreak{}after-\allowbreak{}termination}, \texttt{ICHECK-\allowbreak{}checkpoint-\allowbreak{}on-\allowbreak{}a-\allowbreak{}quarantined-\allowbreak{}session}, \texttt{ICHECK-\allowbreak{}checkpoint-\allowbreak{}over-\allowbreak{}an-\allowbreak{}unverified-\allowbreak{}log}, \texttt{ICHECK-\allowbreak{}leaf-\allowbreak{}hash-\allowbreak{}without-\allowbreak{}its-\allowbreak{}prefix}, \texttt{ICHECK-\allowbreak{}node-\allowbreak{}hash-\allowbreak{}without-\allowbreak{}its-\allowbreak{}prefix}, \texttt{ICHECK-\allowbreak{}leaf-\allowbreak{}and-\allowbreak{}node-\allowbreak{}share-\allowbreak{}a-\allowbreak{}prefix}, \texttt{ICHECK-\allowbreak{}empty-\allowbreak{}tree-\allowbreak{}root-\allowbreak{}not-\allowbreak{}the-\allowbreak{}empty-\allowbreak{}hash}, \texttt{ICHECK-\allowbreak{}subtrees-\allowbreak{}joined-\allowbreak{}in-\allowbreak{}the-\allowbreak{}wrong-\allowbreak{}order}, \texttt{ICHECK-\allowbreak{}last-\allowbreak{}node-\allowbreak{}repeated-\allowbreak{}in-\allowbreak{}an-\allowbreak{}odd-\allowbreak{}tree}, \texttt{ICHECK-\allowbreak{}audit-\allowbreak{}path-\allowbreak{}misplaces-\allowbreak{}the-\allowbreak{}last-\allowbreak{}leaf}\newline \textbf{Rs}: --- (the crate has no CHECKPOINT emission)\newline \textbf{TLA+}: --- (not modelled) \\
I-SIGN-EVENT & --- (keys and services outside the bridge) & --- (outside the reference bridge; no test) & --- (outside the reference bridge; no vector) & \textbf{Py}: --- (no mutant)\newline \textbf{Rs}: --- (not among the sampled rules)\newline \textbf{TLA+}: --- (not modelled) \\
I-SIGN-JWKS & --- (keys and services outside the bridge) & --- (inspection of key sources; no test) & --- (no vector) & \textbf{Py}: --- (no mutant)\newline \textbf{Rs}: --- (not among the sampled rules)\newline \textbf{TLA+}: --- (not modelled) \\
Timer bounds & \texttt{Inv.\allowbreak{}Timers}, \texttt{StepProps[T\_\allowbreak{}total]}, \texttt{StepProps[T\_\allowbreak{}total, T\_\allowbreak{}detach and T\_\allowbreak{}human]}, \texttt{CapEnds} & T30--T30e & V20, V30--V32, V115, V119 & \textbf{Py}: \texttt{bound-\allowbreak{}not-\allowbreak{}enforced}, \texttt{cap-\allowbreak{}bound-\allowbreak{}not-\allowbreak{}enforced}, \texttt{cap-\allowbreak{}forgets-\allowbreak{}closed-\allowbreak{}episodes}, \texttt{cap-\allowbreak{}restarts-\allowbreak{}with-\allowbreak{}a-\allowbreak{}restart}, \texttt{cap-\allowbreak{}counts-\allowbreak{}time-\allowbreak{}resuming}, \texttt{cap-\allowbreak{}disabled}\newline \textbf{Rs}: --- (not among the sampled rules)\newline \textbf{TLA+}: \texttt{Timers: T\_\allowbreak{}detach not stopped in Res}, \texttt{Timers: an expired call keeps its approval}, \texttt{Timers: a session times out early}, \texttt{T\_\allowbreak{}total a disconnect during verification resets the cap}, \texttt{T\_\allowbreak{}total a restart resets the cap}, \texttt{T\_\allowbreak{}total the cap counts time attached}, \texttt{T\_\allowbreak{}total the cap does not end the session} \\
L1 & \texttt{L1} & T26, T26b, T27b, T30c, T128b, T134, T135, T137, T146 & V16--V17, V28, V31--V32, V50, V111b, V115--V118, V120 & \textbf{Py}: \texttt{detach-\allowbreak{}timer-\allowbreak{}restarts-\allowbreak{}on-\allowbreak{}escalation}, \texttt{timer-\allowbreak{}of-\allowbreak{}the-\allowbreak{}tick-\allowbreak{}fires-\allowbreak{}before-\allowbreak{}the-\allowbreak{}request}, \texttt{refused-\allowbreak{}request-\allowbreak{}skips-\allowbreak{}the-\allowbreak{}timers-\allowbreak{}of-\allowbreak{}its-\allowbreak{}tick}, \texttt{expired-\allowbreak{}session-\allowbreak{}served-\allowbreak{}a-\allowbreak{}reattach}, \texttt{cap-\allowbreak{}expiry-\allowbreak{}named-\allowbreak{}as-\allowbreak{}the-\allowbreak{}detachment-\allowbreak{}bound}\newline \textbf{Rs}: \texttt{EXHAUST-\allowbreak{}timers-\allowbreak{}frozen}, \texttt{EXHAUST-\allowbreak{}refused-\allowbreak{}before-\allowbreak{}timers}, \texttt{EXHAUST-\allowbreak{}disconnect-\allowbreak{}refused}\newline \textbf{TLA+}: \texttt{L1 a soft revocation leaves the session in Res} \\
L2 & \texttt{L2} & T6, T26b, T27b, T137 & V11, V17, V30, V31b, V119 & \textbf{Py}: \texttt{human-\allowbreak{}timeout-\allowbreak{}disabled}\newline \textbf{Rs}: --- (liveness is checked by TLC only)\newline \textbf{TLA+}: \texttt{L2 an escalation never times out} \\
Rule Q & \texttt{Inv.\allowbreak{}EndCause}, \texttt{StepProps[rule Q]} & T5b, T22--T22d, T38, T66c, T131, T146, T146b, T156 & V29, V54c, V114, V114b & \textbf{Py}: \texttt{malformed-\allowbreak{}log-\allowbreak{}aborts-\allowbreak{}recovery}, \texttt{operator-\allowbreak{}key-\allowbreak{}added-\allowbreak{}to-\allowbreak{}a-\allowbreak{}quarantined-\allowbreak{}session}\newline \textbf{Rs}: --- (not among the sampled rules)\newline \textbf{TLA+}: \texttt{Q a chain failure terminates the session}, \texttt{Q a recovery failure terminates the session}, \texttt{Q a quarantined session keeps ticking}, \texttt{Q a restart touches a quarantined session} \\
Rule D & \texttt{Inv.\allowbreak{}Ack}, \texttt{StepProps[Rule D, I13]} & T23, T57--T60e, T154 & V12, V33--V37, V38b & \textbf{Py}: \texttt{I2-\allowbreak{}forward-\allowbreak{}policy-\allowbreak{}ignores-\allowbreak{}allowlist}, \texttt{I12-\allowbreak{}unacknowledged-\allowbreak{}call-\allowbreak{}not-\allowbreak{}forwarded-\allowbreak{}again}, \texttt{RuleD-\allowbreak{}drain-\allowbreak{}while-\allowbreak{}detached}, \texttt{session-\allowbreak{}forwarded-\allowbreak{}key-\allowbreak{}shared}, \texttt{RuleD-\allowbreak{}call-\allowbreak{}never-\allowbreak{}forwarded-\allowbreak{}is-\allowbreak{}completed}, \texttt{RuleD-\allowbreak{}directly-\allowbreak{}forwarded-\allowbreak{}call-\allowbreak{}cannot-\allowbreak{}be-\allowbreak{}completed}\newline \textbf{Rs}: \texttt{SESSION-\allowbreak{}forwarded-\allowbreak{}key-\allowbreak{}shared}\newline \textbf{TLA+}: \texttt{Rule D: drain while detached}, \texttt{Ack a call is delivered that was never forwarded}, \texttt{Ack a delivered call is forwarded again}, \texttt{Ack a call is forwarded while an escalation is pending}, \texttt{G1 forward ignores the allowlist}, \texttt{G1 a held call is drained unanswered} \\
Wire format & --- (bytes are outside the model) & T66f, T84, T101--T104, T148--T150, T151, T151b, T155 & \texttt{bytes\_\allowbreak{}kat.\allowbreak{}json}, \texttt{p256\_\allowbreak{}kat.\allowbreak{}json} & \textbf{Py}: \texttt{WIRESERVER-\allowbreak{}attach-\allowbreak{}asks-\allowbreak{}the-\allowbreak{}transport-\allowbreak{}hook}, \texttt{WIRESERVER-\allowbreak{}resolve-\allowbreak{}not-\allowbreak{}authenticated}, \texttt{WIRESERVER-\allowbreak{}extend-\allowbreak{}not-\allowbreak{}authenticated}, \texttt{WIRESERVER-\allowbreak{}attach-\allowbreak{}grant-\allowbreak{}dropped}, \texttt{WIRESERVER-\allowbreak{}extend-\allowbreak{}revert-\allowbreak{}not-\allowbreak{}reported}, \texttt{WIRESERVER-\allowbreak{}operator-\allowbreak{}signature-\allowbreak{}dropped}, \texttt{WIRESERVER-\allowbreak{}not-\allowbreak{}exhausted-\allowbreak{}named-\allowbreak{}as-\allowbreak{}exhausted}, \texttt{WIRESERVER-\allowbreak{}authentication-\allowbreak{}hook-\allowbreak{}not-\allowbreak{}consulted}, \texttt{WIRESERVER-\allowbreak{}revocation-\allowbreak{}request-\allowbreak{}not-\allowbreak{}authenticated}, \texttt{WIRESERVER-\allowbreak{}nesting-\allowbreak{}not-\allowbreak{}bounded}, \texttt{WIRESERVER-\allowbreak{}verified-\allowbreak{}carries-\allowbreak{}event-\allowbreak{}kinds-\allowbreak{}only}, \texttt{WIRESERVER-\allowbreak{}narrowed-\allowbreak{}omits-\allowbreak{}the-\allowbreak{}new-\allowbreak{}version}, \texttt{WIRESERVER-\allowbreak{}tool-\allowbreak{}call-\allowbreak{}response-\allowbreak{}omits-\allowbreak{}the-\allowbreak{}key}, \texttt{WIRESERVER-\allowbreak{}authentication-\allowbreak{}hook-\allowbreak{}gets-\allowbreak{}no-\allowbreak{}message}, \texttt{WIRESERVER-\allowbreak{}disconnect-\allowbreak{}not-\allowbreak{}authenticated}, \texttt{WIRESERVER-\allowbreak{}escalate-\allowbreak{}not-\allowbreak{}authenticated}, \texttt{WIRESERVER-\allowbreak{}submit-\allowbreak{}tool-\allowbreak{}call-\allowbreak{}not-\allowbreak{}authenticated}, \texttt{WIRESERVER-\allowbreak{}answer-\allowbreak{}tool-\allowbreak{}call-\allowbreak{}not-\allowbreak{}authenticated}, \texttt{WIRESERVER-\allowbreak{}delivered-\allowbreak{}not-\allowbreak{}authenticated}, \texttt{WIRESERVER-\allowbreak{}completed-\allowbreak{}not-\allowbreak{}authenticated}, \texttt{WIRESERVER-\allowbreak{}narrow-\allowbreak{}not-\allowbreak{}authenticated}, \texttt{WIRESERVER-\allowbreak{}terminate-\allowbreak{}not-\allowbreak{}authenticated}, \texttt{CANON-\allowbreak{}non-\allowbreak{}ascii-\allowbreak{}escaped}, \texttt{CANON-\allowbreak{}float-\allowbreak{}printed-\allowbreak{}as-\allowbreak{}python}, \texttt{CANON-\allowbreak{}keys-\allowbreak{}by-\allowbreak{}code-\allowbreak{}point}, \texttt{CANON-\allowbreak{}large-\allowbreak{}integer-\allowbreak{}accepted}, \texttt{CANON-\allowbreak{}exponent-\allowbreak{}threshold}, \texttt{canonical-\allowbreak{}wrong-\allowbreak{}order}, \texttt{canonical-\allowbreak{}extra-\allowbreak{}whitespace}, \texttt{canonical-\allowbreak{}ascii-\allowbreak{}escaped}, \texttt{signature-\allowbreak{}wrong-\allowbreak{}scope}, \texttt{resume-\allowbreak{}tag-\allowbreak{}dropped}, \texttt{resume-\allowbreak{}key-\allowbreak{}id-\allowbreak{}dropped}, \texttt{challenge-\allowbreak{}key-\allowbreak{}id-\allowbreak{}dropped}, \texttt{WIRESERVER-\allowbreak{}noncanonical-\allowbreak{}accepted}, \texttt{WIRESERVER-\allowbreak{}add-\allowbreak{}client-\allowbreak{}nonce-\allowbreak{}dropped}, \texttt{WIRESERVER-\allowbreak{}reattach-\allowbreak{}not-\allowbreak{}completed}, \texttt{WIRESERVER-\allowbreak{}challenge-\allowbreak{}key-\allowbreak{}id-\allowbreak{}ignored}, \texttt{WIRESERVER-\allowbreak{}joined-\allowbreak{}carries-\allowbreak{}event-\allowbreak{}kinds-\allowbreak{}only}, \texttt{WIRESERVER-\allowbreak{}reason-\allowbreak{}lost}\newline \textbf{Rs}: \texttt{non-\allowbreak{}canonical-\allowbreak{}accepted}, \texttt{unsigned-\allowbreak{}ops-\allowbreak{}skip-\allowbreak{}auth}, \texttt{auth-\allowbreak{}asked-\allowbreak{}of-\allowbreak{}signed-\allowbreak{}ops}, \texttt{missing-\allowbreak{}field-\allowbreak{}defaults}, \texttt{resume-\allowbreak{}reason-\allowbreak{}hidden}, \texttt{unknown-\allowbreak{}session-\allowbreak{}not-\allowbreak{}named}, \texttt{reattach-\allowbreak{}not-\allowbreak{}completed}, \texttt{join-\allowbreak{}is-\allowbreak{}reattach}, \texttt{answer-\allowbreak{}decision-\allowbreak{}inverted}, \texttt{duplicate-\allowbreak{}call-\allowbreak{}accepted}, \texttt{queue-\allowbreak{}full-\allowbreak{}accepted}, \texttt{resolve-\allowbreak{}not-\allowbreak{}authenticated}, \texttt{attach-\allowbreak{}grant-\allowbreak{}dropped}, \texttt{extend-\allowbreak{}revert-\allowbreak{}not-\allowbreak{}reported}, \texttt{operator-\allowbreak{}signature-\allowbreak{}dropped}, \texttt{not-\allowbreak{}exhausted-\allowbreak{}named-\allowbreak{}as-\allowbreak{}exhausted}, \texttt{disconnect-\allowbreak{}ignores-\allowbreak{}key}, \texttt{authentication-\allowbreak{}hook-\allowbreak{}gets-\allowbreak{}no-\allowbreak{}message}\newline \textbf{TLA+}: --- (not modelled) \\
\bottomrule
\end{longtable}
\endgroup


\subsection{Limitations of the formal model}
\label{app:formal-limitations}

\begin{itemize}[leftmargin=2em]
\item \textbf{Bounded state space.} TLC checks bounded models: two
sessions with an event bound of 2 (one state-growing action per session
after attach) and one session with a bound of 7, both with two client
keys, at most two challenges, two permissions, two tool-call keys and one
restart. Sessions
interact only through the clock tick and the restart. For the safety
properties this removes the number of sessions as a bound. Every state
invariant and step property is quantified over one session at a time; a
step by one session changes only that session's record and its forwarded
calls; and the clock tick and the restart act on each session as a
function of its own record (the restart also chooses which sessions
suffer a recovery fault). Deleting the other sessions from a trace
therefore leaves a trace of the same model with one session and the same
restart bound, and a violation by a session is a violation in that
trace. A one-session model thus covers any number of sessions for the
properties checked as state invariants and step properties. This is an
informal projection argument, not a mechanised one, and it does not
cover the liveness properties L1 and L2. Results do not generalise to
unbounded keys, challenges, permissions, events or restarts without
additional argument.
\item \textbf{Fairness.} L1 and L2 are checked only under weak
fairness of the clock tick, the timeout actions and the completion of
verification (axiom A6), in a one-session model. The event bound also
bounds the $\textsc{Det}\to\textsc{Res}\to\textsc{Det}$ re-attach cycle, so
the model cannot exhibit a client that cycles forever; L1 is verified in
its per-episode form.
\item \textbf{Simplifications of the new rules.} Time does not pass
during a restart in the model, so the cap counts no bridge downtime
(the reference bridge does, and a vector checks it). The detached-call
variants are checked one per configuration, each with one session.
\texttt{TOOL\_CALL\_COMPLETED} changes no state and is not modelled. An
operator-drained call is forwarded at any later point while no escalation
is pending, which over-approximates the reference bridge, which forwards it
at once.
\item \textbf{Abstract cryptography.} The hash chain and the signature
bytes are not modelled (they rely on A3); the prefix check is a verdict flag,
and tampering with the log is a nondeterministic fault. A signature is a
record whose verification is the possession predicate of
Appendix~\ref{app:attacker}, so TLC checks the protocol against an attacker
that holds no client key, not the scheme: unforgeability is assumed, not
model-checked. The cryptographic checks are exercised by the conformance
suite and its mutants (Ed25519 and ECDSA P-256 signatures, a registry of
public keys and scheme tags only).
\item \textbf{The model is not linked to the code by proof.} The model
is written by hand, and no refinement mapping connects it to the reference
implementations. A clean TLC run is evidence about the model. The evidence
about the code is the conformance suite, the shared vectors and the
mutants (Table~\ref{tab:traceability}), which show that the code agrees with
the checked rules on the cases exercised, not that it refines the model.
\item \textbf{No Byzantine faults.} A4 (trusted bridge) and A5
(serialised event emission) are assumed. Faults are limited to log
tampering and recovery failure.
\item \textbf{No multi-hop.} BC $\circ$ BC composition is not modelled.
\item \textbf{Compositional safety is one-directional.} The
compositional safety properties in \S\ref{app:comp-safety} state
invariants over the bridge's state; they do not prove that PACE and
GEC, when composed, preserve these invariants. A full compositional
proof requires stating the corresponding invariants in the PACE and
GEC specifications and showing that the composition preserves them;
this is the target for a future revision.
\end{itemize}

